% TOP_PROBLEM: {"id": 8400027, "title": "O23 — Factoring from Euler's totient", "category_id": 3, "category": {"id": 3, "name": "graph_theory", "display_name": "Graph Theory", "description": "Problems involving graphs, networks, and their properties.", "slug": "graph-theory", "order_index": 3, "created_at": "2026-07-31T15:26:25.671Z"}, "status": "partially_solved", "problem_number": "AMR-083-0027", "set_id": 13}
% FIRST_POSTED: 2026-10-01
% ENTRY_KIND: statement_only
% UnsolvedMath ID 8400027; source catalogue code AMR-083-0027
% !TeX program = lualatex
% Definition and sources only; this file is not an LLM solution attempt.
\documentclass[11pt,a4paper]{article}
\usepackage[margin=25mm]{geometry}
\usepackage{fontspec}
\setmainfont{lmroman10-regular.otf}[BoldFont=lmroman10-bold.otf,
 ItalicFont=lmroman10-italic.otf,BoldItalicFont=lmroman10-bolditalic.otf]
\usepackage{amsmath,amssymb,mathtools}
\usepackage{unicode-math}
\setmathfont{latinmodern-math.otf}
\AtBeginDocument{\let\setminus\smallsetminus}
\usepackage{xurl,hyperref}
\hypersetup{colorlinks=true,urlcolor=blue}
\setcounter{secnumdepth}{0}
\setlength{\parindent}{0pt}
\setlength{\parskip}{5pt}
\setlength{\emergencystretch}{3em}
\begin{document}
\section*{O23 — Factoring from Euler's totient}\label{8400027}
{\small UnsolvedMath ID 8400027 · AMR-083-0027 · Graph Theory · AMR Open Problem Lists}
\subsection{Definitions and mathematical statement}
Euler's totient is the function
\[
\varphi(m)=|\{a\in\{1,\ldots,m\}:\gcd(a,m)=1\}|\qquad(m\ge1).
\]
Assume access to an exact oracle which returns the binary integer $\varphi(m)$ for every positive binary query $m$. The desired output on an input $N\ge2$ is its complete prime factorization $N=\prod_{i=1}^t p_i^{e_i}$, with distinct primes $p_i$ and positive integer exponents $e_i$.

Does a deterministic classical polynomial-time oracle algorithm compute this factorization? Its ordinary bit operations, total number of calls, and query lengths must be bounded by fixed polynomials in the binary length of $N$. A call counts as one oracle step in addition to the cost of constructing its input and reading its answer. The algorithm may make adaptive queries, including queries at integers other than $N$.

Only totient values are returned. In particular, the oracle does not factor either its argument or its output. For a product of two distinct primes, the value of the totient gives useful algebraic information, but the problem includes integers with arbitrarily many distinct prime divisors and arbitrary prime-power multiplicities.

The target is a deterministic reduction on every input. A randomized factoring procedure supplied with the totient, or a deterministic procedure restricted to a special family of integers, does not establish the universal reduction asked for here.
\subsection{Short English statement}
Does exact access to Euler’s totient permit deterministic polynomial-time complete integer factorization?
\subsection{Sources}
\begin{itemize}
\item[S1] ULAM.AI and the UnsolvedMath Contributors, supplied problems.json, record 8400027 (AMR-083-0027), AMR Open Problem Lists. Catalogue identity and source transcription; CC BY 4.0.
\url{https://huggingface.co/datasets/ulamai/UnsolvedMath}.
\item[S2] Adleman \& McCurley - Open Problems in Number Theory Complexity II (1994), O23, p17 / LNCS p307. Original source indexed by AMR-083-0027.
\url{https://mccurley.org/papers/open.ps.gz}.
\item[S3] F. Morain, G. Renault and B. Smith, Deterministic factoring with oracles (2018). Related deterministic reductions and their size restrictions.
\url{https://arxiv.org/abs/1802.08444}.
\end{itemize}
\end{document}
